% FILE: CSE-23_CT-2.tex
\documentclass{article}
\usepackage{amsmath}
\usepackage{amssymb}
\begin{document}

%%%%%%%%%%%%%%%%%%%%%%%%%%%%%%%%%%%%%%%%%%%%%%%%%%
% QUESTION_START
%%%%%%%%%%%%%%%%%%%%%%%%%%%%%%%%%%%%%%%%%%%%%%%%%%
% id=cse23-ct2-secb-q1
% discipline=CSE
% year=2023
% exam_type=CT
% section=B
% ct_number=2
% question_number=1
% topic=Definite Integration
% subtopic=General
% difficulty=Medium
% length=Long
% question_type=Repeated PYQ Pattern
% frequency=1
% appearances=
% OCR_STATUS=VERIFIED
\section*{CT2 (Sec B) - Q1}
Question:
State and derive the relation between the beta and gamma functions.

Solution:
Step 1: State the relationship.
The relation between the Beta function $B(m, n)$ and the Gamma function $\Gamma(x)$ for $m > 0$ and $n > 0$ is given by:
\[
B(m, n) = \frac{\Gamma(m) \Gamma(n)}{\Gamma(m+n)}
\]

Step 2: Express the Gamma functions in exponential form.
By definition, the Gamma function is:
\[
\Gamma(n) = \int_0^\infty e^{-t} t^{n-1} \, dt
\]
Let $t = x^2$. Then $dt = 2x \, dx$. The limits of integration remain $0$ to $\infty$. Substituting these:
\[
\Gamma(n) = \int_0^\infty e^{-x^2} (x^2)^{n-1} (2x \, dx) = 2 \int_0^\infty e^{-x^2} x^{2n-1} \, dx
\]
Similarly, we can write $\Gamma(m)$ using a different dummy variable $y$:
\[
\Gamma(m) = 2 \int_0^\infty e^{-y^2} y^{2m-1} \, dy
\]

Step 3: Multiply the two Gamma functions.
Multiplying the two expressions together:
\[
\Gamma(m) \Gamma(n) = \left( 2 \int_0^\infty e^{-y^2} y^{2m-1} \, dy \right) \left( 2 \int_0^\infty e^{-x^2} x^{2n-1} \, dx \right)
\]
Combining this into a double integral over the first quadrant of the $xy$-plane ($x \ge 0, y \ge 0$):
\[
\Gamma(m) \Gamma(n) = 4 \int_0^\infty \int_0^\infty e^{-(x^2+y^2)} x^{2n-1} y^{2m-1} \, dx \, dy
\]

Step 4: Transform to polar coordinates.
Let $x = r \cos \theta$ and $y = r \sin \theta$.
The differential area element is $dx \, dy = r \, dr \, d\theta$.
The boundary of the first quadrant corresponds to $r$ running from $0$ to $\infty$ and $\theta$ running from $0$ to $\frac{\pi}{2}$.
Substituting these transformations:
\[
\Gamma(m) \Gamma(n) = 4 \int_0^{\pi/2} \int_0^\infty e^{-r^2} (r \cos \theta)^{2n-1} (r \sin \theta)^{2m-1} r \, dr \, d\theta
\]
\[
\Gamma(m) \Gamma(n) = 4 \int_0^{\pi/2} \int_0^\infty e^{-r^2} r^{2n-1+2m-1+1} \cos^{2n-1} \theta \sin^{2m-1} \theta \, dr \, d\theta
\]
\[
\Gamma(m) \Gamma(n) = 4 \int_0^{\pi/2} \int_0^\infty e^{-r^2} r^{2(m+n)-1} \sin^{2m-1} \theta \cos^{2n-1} \theta \, dr \, d\theta
\]

Step 5: Separate the integrals.
Since the limits of integration are independent constants, we can factor the integrals:
\[
\Gamma(m) \Gamma(n) = 2 \left( \int_0^\infty e^{-r^2} r^{2(m+n)-1} \, dr \right) \cdot 2 \left( \int_0^{\pi/2} \sin^{2m-1} \theta \cos^{2n-1} \theta \, d\theta \right)
\]

Step 6: Identify the integral components.
Using the definition established in Step 2:
\[
2 \int_0^\infty e^{-r^2} r^{2(m+n)-1} \, dr = \Gamma(m+n)
\]
Using the trigonometric form of the Beta function:
\[
2 \int_0^{\pi/2} \sin^{2m-1} \theta \cos^{2n-1} \theta \, d\theta = B(m, n)
\]

Step 7: Solve for the Beta-Gamma relation.
Substitute these two identities back into the equation:
\[
\Gamma(m) \Gamma(n) = \Gamma(m+n) B(m, n)
\]
Dividing both sides by $\Gamma(m+n)$:
\[
B(m, n) = \frac{\Gamma(m) \Gamma(n)}{\Gamma(m+n)}
\]

Final Answer:
\[
\boxed{B(m, n) = \frac{\Gamma(m) \Gamma(n)}{\Gamma(m+n)}}
\]
%%%%%%%%%%%%%%%%%%%%%%%%%%%%%%%%%%%%%%%%%%%%%%%%%%
% QUESTION_END
%%%%%%%%%%%%%%%%%%%%%%%%%%%%%%%%%%%%%%%%%%%%%%%%%%

\end{document}